Kerja sama dalam penelitian ini dirancang sebagai kolaborasi keilmuan yang bersifat sukarela dan berbasis saling menghormati.  
Peran mitra ditempatkan sebagai rujukan pengetahuan dan praktik Bahasa Isyarat Indonesia (BISINDO) yang sesuai dengan kaidah penggunaan sehari-hari.  

Bentuk kerja sama yang diharapkan meliputi:  
\begin{enumerate}[noitemsep]
    \item Konsultasi dalam pemilihan kosakata dan bentuk gerakan BISINDO yang sesuai dengan kaidah dan praktik penggunaan sehari-hari.  
    \item Validasi kesesuaian makna dan bentuk gestur yang digunakan dalam penelitian untuk menghindari potensi salah tafsir.  
    \item Pendampingan dalam proses perekaman agar gerakan yang direkam merepresentasikan praktik BISINDO yang benar.  
    \item Keterlibatan penutur atau pengajar Tuli yang direkomendasikan oleh PUSBISINDO sebagai peraga atau model isyarat dalam proses perekaman.  
\end{enumerate}

Keterlibatan penutur atau pengajar Tuli dilakukan atas dasar persetujuan sukarela.  
Peran penutur atau pengajar Tuli diposisikan sebagai mitra keilmuan dan informan ahli, bukan sebagai objek eksperimen.  
Seluruh proses perekaman dan penggunaan data dilaksanakan dengan menjunjung tinggi prinsip etika penelitian dan penghormatan terhadap komunitas Tuli.  